% Obsahový zdroj: import_hodiny/33_Permutace_s_opakovanim.tex a odpovídající PDF (obě strany).
% P1–P5: zdrojové Příklady 1–5; C1–C5: zdrojové Procvičování 1–5.
% Matematický obsah zachován; výklad, odvození a řešení pedagogicky rozvedeny.
% Upřesnění a opravy: viz KMB_29-41_Kontrola.md; původní import zůstává beze změny.
% Rozpis: 4_rocnik/tematicke_plany/Rozpis_vyuky_MAT_4_rocnik.tex, hodina 33.
% Plán: 4_rocnik/tematicke_plany/SPS.pdf, oddíly 2.1–2.3.
% Sazba: sablony/pripravy/TEMPLATE_Priprava.tex a společný priprava.sty.
% Vlastní doplnění: didaktický výklad, zdůvodnění, metodika; žádné nové číslované úlohy.
% Učebnice/09_Kombinatorika nemá obsahovou kapitolu; místní kombinatorické listy ani Realisticky.cz nejsou dostupné.
\documentclass[11pt,a4paper]{article}
% Spolecny styl pro ucitelske pripravy MAT (Work/MAT)
% Verze 1.1, 2026-10-23
\NeedsTeXFormat{LaTeX2e}
\ProvidesPackage{priprava}[2026/10/23 v1.1 Sablona ucitelskych priprav MAT]

% Zakladni typografie a cestina
\RequirePackage[T1]{fontenc}
\RequirePackage[utf8]{inputenc}
\RequirePackage[czech]{babel}
\RequirePackage{lmodern}
\RequirePackage{microtype}
\RequirePackage{amsmath,amssymb}
\RequirePackage{xcolor}
\RequirePackage{enumitem}
\RequirePackage[most]{tcolorbox}
\RequirePackage{tabularx,booktabs,array}
\RequirePackage{fancyhdr}
\RequirePackage{needspace}
\RequirePackage{geometry}
\RequirePackage[hidelinks]{hyperref}

\geometry{a4paper,top=17mm,bottom=18mm,left=23mm,right=23mm}
\definecolor{notegray}{gray}{0.39}
\setlength{\parindent}{0pt}
\setlength{\parskip}{2.5pt plus 1pt minus 0.5pt}
\setlist[itemize]{topsep=2pt,itemsep=1pt,parsep=0pt,leftmargin=1.45em}
\setlist[enumerate]{topsep=2pt,itemsep=2pt,parsep=0pt,leftmargin=1.65em}
\setcounter{secnumdepth}{0}
\renewcommand{\arraystretch}{1.12}

% Zahlavi: vlevo rocnik | globalni tema, vpravo cislo: nazev hodiny
% Zapati zustava zcela prazdne.
\pagestyle{fancy}
\fancyhf{}
\renewcommand{\headrulewidth}{0.35pt}
\renewcommand{\footrulewidth}{0pt}
\setlength{\headheight}{14pt}
\newcommand{\nastavPripravu}[4]{%
  \fancyhead[L]{\small #1\enspace\textbar\enspace #2}%
  \fancyhead[R]{\small #3: #4}%
}

% Minimalisticke cernobile boxy. Kazda definice zvlast.
\tcbset{commonlesson/.style={enhanced,breakable=false,colback=white,colframe=black,
  colbacktitle=white,coltitle=black,boxrule=0.45pt,
  arc=0.8mm,outer arc=0.8mm,left=3mm,right=3mm,top=1.6mm,bottom=1.6mm,
  toptitle=1mm,bottomtitle=0.7mm,fonttitle=\bfseries,
  attach boxed title to top left={xshift=2.5mm,yshift=-1.15mm},
  boxed title style={colback=white,colframe=white,boxrule=0pt,left=1mm,right=1mm}}}
\newtcolorbox{zadaniBox}[1][]{commonlesson,title={Zadání},#1}
\newtcolorbox{prikladBox}[1][]{commonlesson,title={Příklad},#1}
\newtcolorbox{definiceBox}[1][]{commonlesson,title={Definice},#1}

% Spolecny box pro zadani a jeho reseni, bez mezery mezi castmi.
% Pouziti:
%   \begin{prikladReseniBox}
%     Zadani ulohy.
%     \tcblower
%     \begin{align*} ... \end{align*}
%   \end{prikladReseniBox}
% Spodni cast ma jemny sedy podklad a delici linku; lze v ni sazet
% vice radku matematiky. Obe casti maji jen jeden spolecny ramecek.
\tcbset{solutionlesson/.style={
  commonlesson,
  enhanced jigsaw,
  breakable,
  colbacklower=black!3,
  segmentation style={solid,draw=black!25,line width=0.35pt},
  fontlower=\small\color{notegray},
  before lower={\noindent\textbf{Řešení:}\enspace},
  leftlower=3mm,rightlower=3mm
}}
\newtcolorbox{prikladReseniBox}[1][]{solutionlesson,title={Příklad},#1}
\newtcolorbox{zadaniReseniBox}[1][]{solutionlesson,title={Zadání},#1}

% Pomucky pro autora pripravy: sedive poznamky; \reseni pro starsi pripravy
\newcommand{\poznamka}[1]{\par\begingroup\small\color{notegray}\itshape
  \textnormal{\textbf{Poznámka pro učitele.}}\enspace #1\par\endgroup}
\newcommand{\reseni}[1]{\par\begingroup\small\color{notegray}
  \textbf{Očekávané řešení:}\enspace #1\par\endgroup}
\newcommand{\blok}[1]{\needspace{5\baselineskip}\subsection*{#1}}
\newcommand{\vyrok}[1]{\uv{#1}}

\nastavPripravu{4. ročník}{Kombinatorika}{33}{Permutace s opakováním}
\begin{document}


Napiš písmena slova ANANAS. Kdybychom tři kartičky s A tajně očíslovali, dostali bychom více pořadí než při čtení samotných písmen. Ptej se, kolikrát se jedno viditelné slovo objeví v seznamu pořadí očíslovaných kartiček.
\poznamka{Předpoklady: permutace různých prvků, faktoriál a bloková metoda. Cíl: žák spočítá četnosti, odvodí dělitele a vysvětlí, proč záměna stejných kopií netvoří nové pořadí. Nejprve použij P1, pak kontrastní P3 a blokový P4.}

\blok{Výklad a odvození}
\begin{definiceBox}
\textbf{Permutace s opakováním} je uspořádání všech $n$ prvků, z nichž jsou kopie téhož typu nerozlišitelné. Typ $i$ má předem určenou četnost $n_i\in\mathbb N_0$, kde $n_1+\cdots+n_r=n$. Počet různých pořadí značíme $P(n;n_1,\ldots,n_r)$.
\end{definiceBox}

Dočasně všechny kopie očíslujme. Takto rozlišené objekty mají $n!$ pořadí. Po smazání čísel se $n_1!$ výměn kopií prvního typu projeví stejně. Současně lze nezávisle vyměnit kopie druhého typu $n_2!$ způsoby atd. Každé viditelné pořadí má přesně $n_1!\cdots n_r!$ očíslovaných podob. Proto
\[
 P(n;n_1,\ldots,n_r)=\frac{n!}{n_1!\,n_2!\cdots n_r!}.
\]
Dělení je oprávněné právě proto, že počet očíslovaných podob je pro \emph{každé} viditelné pořadí stejný. Typ s četností nula přispívá faktorem $0!=1$. Jsou-li všechny četnosti jedna, vrací se $n!$; je-li jen jeden typ, vychází $n!/n!=1$.

Například u dvou stejných A a jednoho B jsou pouze $AAB,ABA,BAA$. Každé z těchto pořadí se při rozlišení $A_1,A_2$ objeví dvakrát. Toto je malé názorné zdůvodnění dělení faktoriálem četnosti.

U blokových úloh nejprve přepočítáme počet nových objektů a jejich nerozlišitelnost. Blok $AAA$ uvnitř nemá šest různých pořadí: všechna přeuspořádání jsou stejná. Četnosti jsou předem pevné; nepovolujeme volné přidávání dalších kopií znaků.


\poznamka{Průběh: 0–9 min motivace a malý výpis; 9–20 min odvození a P1; 20–31 min P3–P4, P2 a P5 krátce pro srovnání; 31–42 min C1–C3, případně C4; 42–45 min kontrola. C5 je rezerva či domácí práce. Tříbarevné žetony v C4 nepočítej dlouhým vyčíslováním čitatele, nejprve krať.}

\blok{Společně řešené příklady}

% Zdroj P1: import_hodiny/33_Permutace_s_opakovanim.tex, Příklad 1.
\begin{prikladReseniBox}[title={Příklad P1}]
Kolik různých slov (nemusejí dávat smysl) vznikne přeskupením písmen \textsc{ANANAS}?
\tcblower
Je \(6\) písmen: \(A\) třikrát, \(N\) dvakrát, \(S\) jednou. Počet \(\frac{6!}{3!\,2!}=\frac{720}{12}=60\).
\end{prikladReseniBox}


% Zdroj P2: import_hodiny/33_Permutace_s_opakovanim.tex, Příklad 2.
\begin{prikladReseniBox}[title={Příklad P2}]
Kolik pořadí vznikne z písmen \textsc{KOKOS}?
\tcblower
Pět písmen; \(K\) dvakrát, \(O\) dvakrát, \(S\) jednou. \(\frac{5!}{2!\,2!}=30\).
\end{prikladReseniBox}


% Zdroj P3: import_hodiny/33_Permutace_s_opakovanim.tex, Příklad 3.
\begin{prikladReseniBox}[title={Příklad P3}]
Kolik sedmimístných řetězců obsahuje právě čtyři nuly a tři jedničky?
\tcblower
Použijeme všech sedm znaků s pevnými četnostmi: čtyři nuly, tři jedničky. Počet je
\[
 \frac{7!}{4!\,3!}=\frac{7\cdot6\cdot5}{3\cdot2\cdot1}=35.
\]
První znak smí být nula, jde o řetězec. Každé rozložení tří jedniček jednoznačně určuje i všechny čtyři nuly; další volby nezbývají.
\end{prikladReseniBox}


% Zdroj P4: import_hodiny/33_Permutace_s_opakovanim.tex, Příklad 4.
\begin{prikladReseniBox}[title={Příklad P4}]
U slova \textsc{ANANAS} požadujeme, aby všechna tři \(A\) stála vedle sebe.
\tcblower
Spojíme tři A do jediného objektu $AAA$. Řadíme čtyři objekty $AAA,N,N,S$, přičemž dvě N jsou nerozlišitelná. Proto je pořadí $4!/2!=12$. Rozpojením bloku dostaneme právě všechna povolená slova. Uvnitř $AAA$ je pouze jedna viditelná možnost, žádným dalším $3!$ nenásobíme.
\end{prikladReseniBox}


% Zdroj P5: import_hodiny/33_Permutace_s_opakovanim.tex, Příklad 5.
\begin{prikladReseniBox}[title={Příklad P5}]
Kolik různých rozložení pěti shodných červených a dvou shodných modrých žetonů do sedmi míst v řadě?
\tcblower
Použijeme sedm žetonů, přičemž pět červených je navzájem shodných a dva modré také. Dostáváme $7!/(5!\,2!)=7\cdot6/2=21$. Dvě místa obsazená modrými jednoznačně určují celé rozložení.
\end{prikladReseniBox}


\newpage
\blok{Samostatná práce a zásoba k procvičení}

Před vzorcem zapiš seznam typů a jejich četností a zkontroluj jejich součet. U blokové úlohy sestav nový seznam po spojení bloku. C1–C3 jsou základ, C4–C5 zásoba.

% Zdroj C1: import_hodiny/33_Permutace_s_opakovanim.tex, Procvičování, položka 1.
\begin{zadaniReseniBox}[title={Úloha C1}]
Kolik různých řetězců vznikne přeskupením všech písmen \textsc{TATRA}?
\tcblower
Je pět písmen, T dvakrát, A dvakrát a R jednou. Počet je $5!/(2!\,2!\,1!)=30$.
\end{zadaniReseniBox}

% Zdroj C2: import_hodiny/33_Permutace_s_opakovanim.tex, Procvičování, položka 2.
\begin{zadaniReseniBox}[title={Úloha C2}]
Kolik různých řad tvoří tři nuly a dvě jedničky?
\tcblower
Pevné četnosti jsou tři nuly a dvě jedničky, celkem pět míst. Proto $5!/(3!\,2!)=10$ řad.
\end{zadaniReseniBox}

% Zdroj C3: import_hodiny/33_Permutace_s_opakovanim.tex, Procvičování, položka 3.
\begin{zadaniReseniBox}[title={Úloha C3}]
Kolik různých řetězců vznikne přeskupením \textsc{MAMA}, mají-li obě M stát vedle sebe?
\tcblower
Nové objekty jsou $MM,A,A$: tři objekty se dvěma stejnými A. Počet je $3!/2!=3$. Přímo: $MMAA,AMMA,AAMM$.
\end{zadaniReseniBox}

% Zdroj C4: import_hodiny/33_Permutace_s_opakovanim.tex, Procvičování, položka 4.
\begin{zadaniReseniBox}[title={Úloha C4}]
Kolik různých pořadí má osm žetonů: tři červené, tři modré a dva zelené, jsou-li žetony téže barvy nerozlišitelné?
\tcblower
Kontrola četností: $3+3+2=8$. Dělíme výměnami žetonů všech tří barev:
\[
\frac{8!}{3!\,3!\,2!}=\frac{40320}{6\cdot6\cdot2}=560.
\]
\end{zadaniReseniBox}

% Zdroj C5: import_hodiny/33_Permutace_s_opakovanim.tex, Procvičování, položka 5.
\begin{zadaniReseniBox}[title={Úloha C5}]
Kolik různých řetězců vznikne ze slova \textsc{ANANAS}, mají-li obě N stát vedle sebe?
\tcblower
Spojíme N do bloku $NN$. Řadíme pět objektů $NN,A,A,A,S$, z nichž tři A jsou shodná. Dostáváme $5!/3!=20$. Blok $NN$ nemá dvě různá vnitřní pořadí.
\end{zadaniReseniBox}


\poznamka{Typické chyby: dělení četností místo jejich faktoriálů; násobení pořadím uvnitř nerozlišitelného bloku; záměna délky slova a počtu různých písmen. U ANANAS je $n=6$, počet typů je tři. Příště připustíme libovolnou volbu znaku na každé pozici a porovnáme ji s dnešními pevnými četnostmi.}
\end{document}
